% Expose IX, sections 11.0--11.2.
% French transcription lines 2956--3027; authority scan physical pp. 460--463
% (scan indices 459--462; source folios 448--451). Physical p. 463 was
% checked through the exact 11.3 boundary. All pages were compared in four
% overlapping 1100-dpi bands; the text and formulas are clear at that
% resolution, so no higher-resolution escalation was necessary.

\medskip
\noindent
11. \underline{Connected components of the N\'eron model: duality
relation and asymptotic behaviour}

\smallskip
\noindent
11.0. The purpose of the present section is to study the finite \'etale
$k$-group

\begin{equation}\tag{11.0.1}
\Phi_o=A_o/A_o^{,o}
\end{equation}
already considered in section 1. For every prime number $\ell$, denote
by

\begin{equation*}
\Phi_o(\ell)
\end{equation*}
its $\ell$-primary component. Thus knowledge of $\Phi_o$ is equivalent
to knowledge of the $\Phi_o(\ell)$, by the formula

\begin{equation}\tag{11.0.2}
\Phi_o=\coprod_\ell\Phi_o(\ell).
\end{equation}

Notice that the groups $\Phi_o$ and $\Phi_o(\ell)$ are determined,
respectively, by the ordinary finite groups $\Phi_o(\bar k)$ and
$\Phi_o(\ell)(\bar k)$ together with the natural action of

\begin{equation*}
\pi_o=\operatorname{Gal}(\bar k/k)
\end{equation*}
on them. On the other hand, it is clear that $\Phi_o(\bar k)$ is also
the group of connected components of the N\'eron model of
$A_{\widetilde K}$, where $\widetilde K$ is the fraction field of the
strict henselization of $S$ relative to the chosen separable closure
$\bar k$ of $k$. Thus, in studying $\Phi_o$, we may reduce to the case
where $S$ is strictly local; the actions of

\begin{equation*}
\pi_o\simeq\operatorname{Gal}(\widetilde K/K)
\end{equation*}
on $\Phi_o(\bar k)$ are then recovered automatically by transport of
structure. We could likewise reduce further to the case where $S$ is
complete, using the base change $\widehat S\longrightarrow S$ and the
fact that formation of the N\'eron model commutes with it
[21, Th.~3.4 b)].

\smallskip
\noindent
11.1. Suppose, then, that $S$ is strictly local, so that $\Phi_o$ may be
identified with the ordinary finite group $\Phi_o(k)$. We obtain

\begin{equation}\tag{11.1.1}
\Phi_o(k)\simeq A_o(k)/A_o^o(k)
 \simeq A(S)/A^o(S)
 \simeq A(K)/A(K)^o,
\end{equation}
where the first equality follows from the fact that $A_o$ is smooth over
$k$, the second from Hensel's lemma since $A$ is smooth over $S$ and
$S$ is henselian, and the third from the relation
$A(S)\xrightarrow{\sim}A(K)$ and from the convention

\begin{equation}\tag{11.1.2}
A(K)^o\mathrel{\mathop:}=A^o(S).
\end{equation}
Thus $A(K)^o$ appears as a natural subgroup of finite index in $A(K)$
(and is not to be confused with $A^o(K)=A(K)$). The isomorphisms
(11.1.1) are plainly functorial in $A$ and compatible with transport of
structure in $(S,A_K)$, and in particular with the action of any Galois
groups.

Now let $\ell$ be a prime number different from $p$. Then $A^o$ is
$\ell$-divisible; more precisely, $\ell\operatorname{id}_{A^o}$ is an
\'etale surjective morphism. It consequently induces a surjective
endomorphism of $A^o(S)=A(K)^o$, which is therefore an
\underline{$\ell$-divisible group}. It follows at once that, for every
integer $n\geq 0$ prime to $p=\operatorname{char}(k)$,

\begin{equation}\tag{11.1.3}
{}_n\Phi_o(k)\simeq{}_nA(K)/{}_nA(K)^o,
 \qquad (n,p)=1.
\end{equation}
By the definition (2.2.3) of the fixed part of ${}_nA_K$, we have

\begin{equation}\tag{11.1.4}
{}_nA(K)^o=({}_nA_K)^f(\bar K)
 \simeq T_n(A_K)^f(\bar K)
 \otimes_{\mathbb Z_n}\mathbb Z/n\mathbb Z.
\end{equation}
Taking $n$ to be a power of $\ell$, and putting as in (2.5.3)

\begin{equation}\tag{11.1.5}
U=T_\ell\bigl(A(\bar K)\bigr),
 \qquad
T_\ell\bigl(A(\bar K)\bigr)^f=U^I,
\end{equation}
where $U^I$ denotes the submodule of $U$ consisting of invariants under
the inertia group $I$, we may write (11.1.3) in the form

\begin{equation}\tag{11.1.6}
{}_n\Phi_o(k)\simeq (U_n)^I/(U^I)_n,
 \qquad n=\ell^{r+1},
\end{equation}
where the subscript $n$ denotes reduction modulo $n$, that is, tensor
product over $\mathbb Z_\ell$ with
$\mathbb Z_\ell/n\mathbb Z_\ell$. The transition morphisms between the
groups occurring on the two sides of (11.1.6) being evident, passage to
the inductive limit over the powers of $\ell$ gives the following
result.

\medskip
\noindent
\underline{Proposition 11.2.} \underline{Suppose that $S$ is strictly
local, and let $\ell$ be a prime number different from $p$. Then there is
a canonical isomorphism}

\begin{equation}\tag{11.2.1}
\Phi_o(\ell)\simeq
 \frac{(U\otimes D_\ell)^I}{U^I\otimes D_\ell},
\end{equation}
\underline{where $U$ is defined in (11.1.5) and
$D_\ell=\mathbb Q_\ell/\mathbb Z_\ell$.}

Thus the group $\Phi_o(\ell)$ can be described solely in terms of the
Tate module $U=T_\ell(A(\bar K))$, regarded as a Galois module under the
action of the inertia group $I$.
